\section{The hereditary forest estimate}\label{app:forest}

This appendix explains why the local deficit inequality applies to an
arbitrary covering word. A shortest word gives a minimum overlap path,
but that path need not visit every rotation class consecutively. The
forest construction replaces it by auxiliary complete-block paths and
records the cost of the replacement. The classes in different components
partition all the rotation classes. A block borrowed at the entrance to
one component is required to be new only to that component.

The construction uses a particular minimum path. Call a minimum path
\emph{reduced} if it has no edge from $u$ to its second cyclic rotation
$\sigma^2(u)$. Such a path exists. Among the minimum paths, choose one
with as many weight-one edges as possible. If it contained
$u\to\sigma^2(u)$, remove $\sigma(u)$ from its other location and
insert it between these two vertices. Removal cannot increase the weight,
by the triangle inequality, and insertion replaces one weight-two edge
by two weight-one edges. Moreover, removal deletes no weight-one edge.
The only possible weight-one predecessor of $\sigma(u)$ is $u$, whose
successor is already $\sigma^2(u)$, and the only possible weight-one
successor is $\sigma^2(u)$, whose predecessor is already $u$.
The new path would therefore still be minimum and have more weight-one
edges, a contradiction.

We state the properties of the forest construction for this reduced
minimum path. The construction and its residual inequality are inherited
from the lower-bound framework~\cite{pantone-framework}. The summation
argument after the statement is included here.

\begin{lemma}[Forest decomposition]\label{lem:forest-data}
Put $K=n+1\ge8$ and $T=n!$. A reduced minimum overlap path determines a finite
set of complete-block path components with a distinguished root. Each
nonroot component is classified as cheap or expensive. For component
$v$, let $b_v$ be its number of blocks, $q_v$ its number of runs,
$D_v$ its total deficit, $m_v$ its number of light segments, and
$d_v^-,d_v^+$ its first and last deficits. There are nonnegative
integers $A_v$ and $y_v$ such that
\[
 \sum_v b_v=T,\qquad b_v=nq_v-D_v,\qquad
 z_v=q_v+A_v,\qquad A_v\ge m_v-1.
\]
For the root and cheap components, $y_v=z_v$; for expensive components,
$y_v\ge z_v+1$.

For a cheap component with at least two runs, one can prepend an
entrance block to obtain another class-distinct path. Its block count
increases by one, its number of runs and $z_v$ are unchanged, and its
first deficit decreases by one. Its last deficit is unchanged.
A cheap component with one run has $y_v=1$ and
$d_v^-=d_v^+\ge1$.

There is a directed absorption relation on the components. Each component
has at most one incoming and one outgoing edge, and every target is a
cheap nonroot component. Every edge is either an isolated loop at a cheap
one-run component of deficit one, or an ordinary edge $v\to w$ with
\[
 d_v^+\ge2,\qquad d_w^-\ge2,\qquad d_v^++d_w^-\ge n+1.
\]
If $V$ is the number of cheap nonroot components with no incoming edge,
and $R=\sum_v y_v$, then
\begin{equation}\label{eq:residual}
 S(K)\ge K+(K+1)T-3+R+V.
\end{equation}
\end{lemma}

The last inequality accounts for the original path as well as the
auxiliary components. In particular, it does not assume that entrance
blocks borrowed by different components are distinct. This is the reason
for proving the local inequality on every genuine segment, including
consecutive factors of an entrance-extended path.

\begin{theorem}[Hereditary forest estimate]\label{thm:forest}
Suppose $0<\delta<1$ and $1\le\beta\le\alpha:=n-\delta$.
If every genuine light segment with $r$ runs and total deficit $D$
satisfies $D\ge\delta r-\beta$, then
\[
 S(K)\ge K+(K+1)n!-3+
 \ceil{\frac{n!-\beta}{n-\delta}}.
\]
\end{theorem}

\begin{proof}
The first step recovers endpoint information from the hypothesis.
Write $F(d)=\max(0,d-\delta)$. Deleting the first deficit from a
segment and applying the hypothesis to the remainder gives an extra
term $d^- -\delta$. Applying the hypothesis without deletion gives
an extra term zero. Taking the larger of these two estimates gives
\[
 D\ge\delta r-\beta+F(d^-).
\]
There is an identical estimate at the last endpoint. When $r\ge2$,
we can delete either endpoint, both, or neither, and obtain
\begin{equation}\label{eq:two-endpoints}
 D\ge\delta r-\beta+F(d^-)+F(d^+).
\end{equation}
The empty remainder causes no difficulty because $0\ge-\beta$.
The two endpoints in \eqref{eq:two-endpoints} must be distinct;
we shall handle one-run components separately.

Apply Lemma~\ref{lem:forest-data} and put
$s_v=b_v-\alpha y_v$. Summing the segment estimates within one
component gives
\[
 b_v-\alpha z_v
 \le m_v\beta-\alpha A_v-F(d_v^+)
 \le\beta-F(d_v^+),
\]
because $A_v\ge m_v-1$ and $\alpha\ge\beta$. Thus
\[
 s_v\le\beta-F(d_v^+)\quad\hbox{at the root},\qquad
 s_v\le\beta-\alpha-F(d_v^+)\quad\hbox{at an expensive component}.
\]
For a cheap component with at least two runs, apply the same argument
to its entrance-extended path and use both endpoint estimates. Define
$A(d)=\max(0,d-1-\delta)$ and $c=\beta-1$. The additional block gives
\[
 s_v\le c-A(d_v^-)-F(d_v^+).
\]
If the two ends lie in different light segments, apply the one-end
estimate to each of those segments. If they lie in the same segment,
that segment has at least two runs and \eqref{eq:two-endpoints} applies.
For a cheap component with a single run, the exact identity is
\[
 s_v=(n-d_v^-)-(n-\delta)=\delta-d_v^-=-F(d_v^-).
\]

Put $\widehat A(d)=\min(A(d),c)$. For a cheap component, let
$i_v,o_v\in\{0,1\}$ indicate whether it has an incoming and an
outgoing absorption edge. Then both the multi-run and single-run
estimates imply
\[
 s_v\le c-i_v\widehat A(d_v^-)-o_vF(d_v^+).
\]
For a single run, this follows from $0\le\widehat A\le c$.
Clipping the incoming charge is what permits one formula to cover that
case without counting its only deficit twice.

Assign an intercept $c_v$ to each component: use $\beta$ at the root,
$c$ at cheap nonroots, and $\beta-\alpha$ at expensive nonroots.
Assign the incoming charge $a_v=\widehat A(d_v^-)$ at cheap vertices
and zero elsewhere, and the outgoing charge $b'_v=F(d_v^+)$.
There is one exception. At a vertex carrying an isolated special loop,
use $c_v=\delta-1$ and $a_v=b'_v=0$; its surplus is exactly
$\delta-1$.

The preceding bounds give
\[
 s_v+i_va_v+o_vb'_v\le c_v.
\]
For an ordinary edge $v\to w$, its two deficits give
\[
 A(d_w^-)+F(d_v^+)\ge n-2\delta,\qquad
 F(d_v^+)\ge2-\delta.
\]
Consequently its charges cover the target intercept, since
\[
 b'_v+a_w-c_w
 =\min\{F(d_v^+)+A(d_w^-)-c,\ F(d_v^+)\}
 \ge\min\{n+1-\beta-2\delta,\ 2-\delta\}>0.
\]
For a special loop, the corresponding excess is $1-\delta>0$.

Now sum the vertex inequalities and the edge inequalities. Each incoming
and outgoing charge occurs exactly once on either side and cancels.
The remaining intercepts belong to vertices with no incoming edge.
These are the root, the $V$ unabsorbed cheap vertices, and all the
expensive vertices. The latter have nonpositive intercept. It follows that
\[
 T-\alpha R=\sum_v s_v\le\beta+(\beta-1)V.
\]
Since $\beta\le\alpha$, this implies
$T\le\alpha(R+V)+\beta$. The number $R+V$ is an integer, so
\eqref{eq:residual} gives the asserted ceiling bound.
\end{proof}
